% discussion.tex -- Discussion
\section{Discussion}

Two results organize this study. A persistent change in how traders
responded to their neighbors was sufficient to move the modeled market
from mixed behavior to a Sell majority, and what that majority did to
the price was decided at execution rather than on the lattice. The
altered rule remained in place when a targeted trader left Sell,
allowing the new response to act whenever its neighborhood input
recurred, whereas a matched one-time Sell pulse changed only the
current state and produced no majority crossings in the tested 50-seed
controls. Persistent Rule~19286 activation produced crossings in every
seed at the two larger fractions, and several other single-entry
changes also produced majorities, although most screened rules did not.
Collective selling therefore depended on the response being changed as
well as on the fraction of traders adopting it. Because paired paths
shared the same realized market state, their later divergence followed
the intervention by construction. The measured transition curve
supplied the additional result: it described how often that local
change spread to a majority and how the outcome depended on the
activation fraction.

The local rule dynamics could sustain the transition without a price
signal. Holding boundary traders inactive left the interior free to
evolve under neighbor interactions alone, and Sell majorities still
developed. In the displayed crossing trajectories, Sell-rich regions
formed, enlarged, and merged, which connected the local responses to
the aggregate Sell fraction in the way that evolving spatial patterns
do in other nonequilibrium lattice
systems~\cite{hinrichsen2000nonequilibrium,bray1994theory}. The
dependence on activation fraction resembled the spread found in
threshold-contagion and bootstrap-percolation
models~\cite{granovetter1978threshold,watts2002simple,aizenman1988metastability,holroyd2003sharp},
where a seeded fraction likewise separates local from global spread.
Because the majority grew out of finite domains, its location depended
on the size of the lattice and on the time allowed. The response
narrowed between the smallest and largest lattices on the refined
grids, yet its midpoint continued to fall as the observation window
increased, so the measured activation fraction was a property of both
size and time. Establishing an asymptotic transition would require
horizon convergence and a specified scaling
model~\cite{binder1981finite,landau2021guide}. The spatial evidence
also stops short of classical nucleation, since no free-energy barrier,
activation rate, surface tension, or critical droplet was identified.
The observed mechanism was growth and merger of Sell-rich domains over
finite times.

A Sell state generated a desired order, but trader and dealer resources
determined whether that order could move the price. Trader holdings
first limited order eligibility, opposing eligible orders then cleared
against one another, and the dealer absorbed only the unmatched flow
that remained, up to its cash and inventory bounds. Those bounds
connect this model to inventory-based accounts of
intermediation~\cite{garman1976market,stoll1978supply,ho1981optimal,braunmunzinger2018agent},
and only the imbalance in filled orders entered the price-impact
relation~\cite{kyle1985continuous,farmer2005predictive,bouchaud2009how}.
At zero dealer capacity, balanced execution left the quote fixed,
whereas unlimited capacity allowed every eligible order to execute. The
same illustrated seed consequently reached a Sell majority with either
a 0.40\% or a 25.8\% drawdown. At 18\% activation with baseline cash,
tenfold trader shares produced an 83.5\% median drawdown against the
unlimited dealer, but finite capacity limited the corresponding
drawdowns to 3.92\% or 0.399\%. More shares made additional selling
eligible, but additional unmatched selling lowered the quote only while
the dealer could continue buying. A lattice full of sellers therefore
lowered the price only as far as someone was able to buy.

Price feedback allowed the execution constraint to influence the spread
of selling as well as its financial consequence. Near the activation
window, capacity $h=0.1$ raised majority probability by 38 percentage
points relative to the unlimited dealer. The quote determined boundary
trader states through a moving center and an adaptive decision width,
and freezing each path's width at its pre-switch value removed the
low-capacity advantage in the point estimates, linking that advantage
to adaptation after the switch. Frozen widths still differed across
seeds and capacities, the moving center and financial accounts
continued to evolve, and other non-monotone capacity differences
remained, so width adaptation did not explain the entire capacity
response. The impact sweep placed a further bound on the price
interpretation: drawdowns remained ordered in capacity across a decade
of impact coefficients, but their magnitudes changed by a factor of
seven to ten. Thus, limiting execution could reduce drawdown while
price feedback increased the probability of a Sell majority. The two
outcomes measured different consequences of the same resource
constraint, and a single crash indicator would have combined them.

Several modeling choices limit how far these conclusions extend.
Traders updated synchronously on a regular lattice using one baseline
rule, and the totalistic score assigned the same response to distinct
neighbor compositions, so the rule screen did not test arbitrary
composition-based responses. Unit orders settled at a common quote,
and the dealer had no strategic quoting policy, risk penalty, or source
of replenishment beyond reversing flow. These assumptions made the
financial constraints explicit but supplied neither an order book nor
an empirical calibration of market depth, price impact, or time.
Ten no-switch seeds under each boundary remained below a Sell majority
for 100,000 additional steps, which extended the observed baseline
persistence without proving global stability, and neither the
1,000-step warm-up nor the final 500-step summary established a
stationary distribution. Asynchronous updates, other interaction rules,
and an order book would allow the same question to be examined under
different assumptions: whether a persistent local response can spread
through the market, and how available resources limit the resulting
price decline.
